\documentclass[10pt,a4paper]{amsart}
\usepackage[margin=19mm]{geometry}
\usepackage{amsmath,amssymb,mathtools}
\usepackage[scr=rsfs,cal=euler]{mathalfa}
\usepackage{xfrac}
\usepackage[unicode,hidelinks,bookmarks=false]{hyperref}
\newcommand{\ie}{i.e.\ }
\newcommand{\cf}{cf.\ }
\newcommand{\resp}{respectively\ }
\newcommand{\Subsection}{\subsection}
\providecommand{\nobreakdash}{\nobreak\hbox{-}}
\setlength{\emergencystretch}{2em}
\multlinegap0pt
\usepackage{bm}
\usepackage{bbm}
\usepackage{dsfont}
\usepackage{xspace}
\usepackage{cancel}
\usepackage{mathtools}
\usepackage{cleveref}
\usepackage{amsthm}
\makeatletter
\@ifundefined{thm}{\newtheorem{thm}{Thm}}{}
\@ifundefined{lemma}{\newtheorem{lemma}[thm]{Lemma}}{}
\@ifundefined{prop}{\newtheorem{prop}[thm]{Proposition}}{}
\makeatother
\newcommand{\F}{\mathbb{F}}
\newcommand{\bv}{\bm{v}}
\title[Random Gabidulin Codes Achieve List Decoding Capacity in the]{Random Gabidulin Codes Achieve List Decoding Capacity in the Rank Metric}
\author{Zeyu Guo, Chaoping Xing, Chen Yuan, Zihan Zhang}
\date{Source: arXiv:2404.13230v2 (CC BY 4.0)}
\begin{document}
\maketitle

\noindent\emph{Source and licence.} Random Gabidulin Codes Achieve List Decoding Capacity in the Rank Metric. Version arXiv:2404.13230v2 (CC BY 4.0), distributed under the Creative Commons Attribution 4.0 International licence, as recorded in the arXiv licence metadata for this exact version and shown on the abstract page https://arxiv.org/abs/2404.13230v2. Full rights notice: licenses/arxiv-2404.13230-CC-BY-4.0.txt.\par\smallskip

\noindent\textit{Editorial scope.} Counted unit: the Lemma whose source label is \texttt{order-l-cha} in the article (source0.tex lines 1406--1422, source label \texttt{order-l-cha}), delivered together with the article's own complete original proof of it (source0.tex lines 1424--1539, one \texttt{proof} environment of 116 lines). No mathematics is written, repaired or re-derived: statement and proof are the article's own text line for line. Required context retained uncounted: lm:dim (lines 819--824); prop:equivalent (lines 1294--1307). Every definition, display and label of those lines is retained unchanged. Omitted from this package and not redistributed: 1--818, 825--1293, 1308--1405, 1423--1423, 1540--2792. Nothing inside a retained excerpt is omitted or altered: every retained body is delivered exactly as the article's lines are written, so no omission receipt of any kind accompanies this package. Citations: the retained excerpts carry no citation command at all, so no bibliography entry of the article is transcribed and no reference is redistributed. Reference adaptation: none was needed and none was made; every reference inside the delivered unit resolves inside it, so no unresolved reference or citation is produced. Printed slips kept literal: no printed word, symbol, hypothesis or number of the counted statement or of its proof was altered. Layout and class: the article's own class is \texttt{article} with packages including geometry, amsmath, amssymb, amsfonts, amsthm, bm, bbm, mathtools, mathrsfs, enumitem, setspace, inputenc, algpseudocode, algorithmicx; the wrapper is the lane's pinned \texttt{amsart} editorial wrapper at 10pt on A4, which restates the theorem-like environments the retained text needs and adds the article's own macro definitions that the retained text needs (\texttt{\textbackslash F}, \texttt{\textbackslash bv}), with extra wrapper packages (bm, bbm, dsfont, xspace, cancel, mathtools, cleveref). The delivered numbering is the unit's own. Assets: no retained span contains a figure environment, a graphics-inclusion command, a PDF-page inclusion, a table or a TikZ picture; no image file is copied into this package, the package declares no asset dependency and no image file is redistributed with it. Licence: version arXiv:2404.13230v2 of this article is distributed under the Creative Commons Attribution 4.0 International licence, as recorded in the arXiv licence metadata for this exact version and shown on the abstract page https://arxiv.org/abs/2404.13230v2. A full rights notice is delivered at licenses/arxiv-2404.13230-CC-BY-4.0.txt. Characters: every retained excerpt is pure ASCII, so the delivered engine, XeLaTeX with its default Latin Modern text font, reports no missing character.
\begin{lemma}\label{lm:dim}
Let $A_1,\ldots,A_m,B_1,\ldots,B_m$ be subspaces of $\F^n$. Then 
\[
\dim\left(\bigcap_{i=1}^m (A_i+B_i)\right)\leq \dim\left(\bigcap_{i=1}^m A_i\right)+\sum_{i=1}^m \dim(B_i).
\]
\end{lemma}

\begin{prop}\label{prop:equivalent}
Let $n$ and $k$ be integers with $n\geq k\geq 0$. Let $V_1,\ldots,V_\ell$ be subspaces of $\F_q^n$, each of dimension at most $k$. Then the following are equivalent. 
%\ZGnote{We need to assume $V_1,\dots,V_\ell$ are distinct, right? Otherwise the kernel pattern may have order less than $\ell$. And could $V_i$ be the zero subspace?} 
\begin{enumerate}
\item\label{it1} There exist integers $\delta_1,\ldots,\delta_\ell\geq 0$ such that for every nonempty $\Omega \subseteq [\ell]$, 
\begin{equation}\label{eqn:copypattern}
\dim\left(\bigcap_{i\in \Omega} V_i\right)\leq k-\sum_{i\in \Omega}\delta_i.
\end{equation}
\item\label{it2} The pattern $(T_1,\ldots,T_k)$ that consists of $\delta_i$ copies of $V_i$ for $i\in [\ell]$ and additional $k-\sum_{i=1}^\ell \delta_i$ copies of $\{0\}$ is a generic kernel pattern.\footnote{We assume that the set of copies of each $V_i$ is disjoint from both the set of copies of any other $V_{i'}$ (even if $V_i = V_{i'}$) and the set of the additional $k - \sum_{j=1}^{\ell} \delta_{j}$ copies of $\{0\}$ (even if $V_i = \{0\}$).} That is, for every nonempty $\Omega'\subseteq [k]$, 
\begin{equation}\label{eqn:pattern}
\dim\left(\bigcap_{i\in \Omega'} T_i\right)\leq k-|\Omega'|.
\end{equation}
\end{enumerate}
\end{prop}

\begin{lemma}\label{order-l-cha}
Let $n$, $k$, and $d$ be integers with $n\ge k\geq d\geq 0$.
Let $V_1,\ldots,V_\ell$ be subspaces of $\F_q^n$, each of dimension at most $k$. Then the following are equivalent. 
\begin{enumerate}
\item\label{eqiv:1} There exists a generic kernel pattern $\mathcal{T}:=(T_1,T_2,\dots,T_k)$
%, where each $T_i$ is a subspace of $\F_q^n$, 
consisting solely of copies of $V_1,\ldots,V_\ell$ and an additional $d$ copies of $\{0\}$.
\item\label{eqiv:2} There exist integers $\delta_1,\ldots,\delta_\ell\geq 0$ such that $\sum_{i=1}^{\ell}\delta_i=k-d$ and for every nonempty $\Omega\subseteq [\ell]$, 
\begin{equation}\label{eq:cond2}
\dim\left(\bigcap_{i\in \Omega}V_i\right)\leq k-\sum_{i\in \Omega}\delta_i.
\end{equation}
\item \label{eqiv:3} For all partitions $P_1\sqcup P_2\sqcup\cdots\sqcup P_s=[\ell]$, we have 
\begin{equation}\label{eq:cond3}
\sum_{i=1}^s\dim\left(\bigcap_{j\in P_i} V_j\right)\leq (s-1)k+d.
\end{equation}
\end{enumerate}
\end{lemma}

\begin{proof}
By \cref{prop:equivalent}, \cref{eqiv:1} is equivalent to \cref{eqiv:2}. 
And \cref{eqiv:2} implies \cref{eqiv:3} because for any partition $P_1\sqcup P_2\sqcup\cdots\sqcup P_s=[\ell]$, we have 
\[
\sum_{i=1}^s \dim\left(\bigcap_{j\in P_i} V_j\right)\stackrel{\eqref{eq:cond2}}{\leq} \sum_{i=1}^s \left(k-\sum_{j\in P_i}\delta_j\right)=sk-(k-d)=(s-1)k+d.
\]
Finally, we prove that \cref{eqiv:3} implies \cref{eqiv:2} via an induction on $\ell$. 
When $\ell=1$, \cref{eqiv:2} and \cref{eqiv:3} are the same.
Now consider $\ell\geq 2$ and assume that \cref{eqiv:3} implies \cref{eqiv:2} for $\ell'<\ell$.
We say a partition $P_1\sqcup P_2\sqcup\cdots\sqcup P_s$ of $[\ell]$ is \emph{tight} if \eqref{eq:cond3} holds with equality.

Suppose \cref{eqiv:3} holds. We claim that there exist subspaces $V'_i\supseteq V_i$ of dimension at most $k$,
%of dimension at most $k$, 
$i=1,\dots,\ell$, such that for these new subspaces, which we call \emph{padded subspaces}, \cref{eqiv:3} still holds and there exists a partition $P_1\sqcup P_2\sqcup\cdots\sqcup P_s$ of $[\ell]$ with $s\geq 2$ that is tight.

The remaining proof consists of two parts: We will first prove the above claim. Then we will show that \cref{eqiv:2} holds for the padded subspaces $V_1',\dots,V_\ell'$.
This suffices since \cref{eqiv:2} then holds for the original subspaces $V_1,\dots,V_\ell$ as well by the fact that $\bigcap_{i\in \Omega}V_i \subseteq  \bigcap_{i\in \Omega}V_i'$ for all nonempty $\Omega\subseteq [\ell]$.

For convenience, define $V_\Omega=\bigcap_{i\in\Omega} V_i$ for nonempty $\Omega\subseteq [\ell]$.
Applying \cref{eqiv:3} to the coarsest partition of $[\ell]$ shows that $\dim V_{[\ell]}\leq d$.

To prove the claim, we repeatedly select pairs $(V_i, \bv)$, where $\dim V_i < k$ and $\bv \in \mathbb{F}_q^n \setminus V_i$, and then replace $V_i$ with $V_i \oplus \mathrm{span}(\bv)$, until one of the partitions of $[\ell]$ becomes tight.\footnote{Note that at least one partition becomes tight when (or before) $\dim V_i=k$ for all $i\in [k]$. This is because (1) if all $V_i$ have dimension $k$, then we would have $\sum_{i=1}^\ell \dim V_i=\ell k\geq (\ell-1)k+d$, and (2) replacing $V_i$ with $V_i \oplus \mathrm{span}(\bv)$ increases LHS of \eqref{eq:cond3} by at most one.}
If the tight partition $P_1\sqcup \cdots\sqcup P_s$ arising this way satisfies $s\geq 2$, then we are done. So assume $s=1$, i.e., $\dim V_{[\ell]}=d$. In this case, we show that we can continue padding the subspaces $V_i$ until another partition with $s\ge2$ becomes tight, while maintaining $\dim V_{[\ell]} = d$.
%For $i\in [\ell]$, let $\widetilde{V}_i$ be a complement of $V_{[\ell]}$ in $V_i$.

For $i\in [\ell]$, applying \cref{eqiv:3} to the partition $\{i\}\sqcup ([\ell]\setminus \{i\})$ of $[\ell]$ shows
\begin{equation}\label{eq:special-partition}
\dim V_i+\dim\left(\bigcap_{j\in[\ell]\setminus\{i\}}V_j\right)\leq k+d.
\end{equation}
By inclusion-exclusion and the fact that $\dim V_{[\ell]}=d$, we have that for $i\in [\ell]$,
\[\dim\left( V_i+\bigcap_{j\in[\ell]\setminus\{i\}}V_j\right)=\dim V_i+\dim\left(\bigcap_{j\in[\ell]\setminus\{i\}}V_j\right)-\dim V_{[\ell]}\stackrel{\eqref{eq:special-partition}}{\leq} k.
\]
If there exists $i\in[\ell]$ such that $\dim\left( V_i+\bigcap_{j\in[\ell]\setminus\{i\}}V_j\right)=k$, then
\eqref{eq:special-partition} must hold with equality. In this case, the partition $\{i\}\sqcup[\ell]\setminus\{i\}$ is tight and we are done.
So assume
\begin{equation}\label{eq:less-than-k}
\dim\left( V_i+\bigcap_{j\in[\ell]\setminus\{i\}}V_j\right)<k\leq n \quad \text{ for all } i\in[\ell].
\end{equation}
Fix arbitrary $i\in[\ell]$. By \eqref{eq:less-than-k}, there exists a vector $\bv\in\F_q^n\setminus (V_{i}+\bigcap_{j\in[\ell]\setminus\{{i}\}}V_j)$. Let $V_i'=V_i \oplus \mathrm{span}(\bv)$.
Then
\begin{align*} &\dim\left(V_i'\cap\left(\bigcap_{j\in[\ell]\setminus\{i\}}V_j\right)\right)\\
&=\dim V_i'+\dim\left(\bigcap_{j\in[\ell]\setminus\{i\}}V_j\right)-\dim\left(V_i'+\left(\bigcap_{j\in[\ell]\setminus\{i\}}V_j\right)\right)\\
&=\dim V_i+1+\dim\left(\bigcap_{j\in[\ell]\setminus\{i\}}V_j\right)-\dim\left(V_i+\left(\bigcap_{j\in[\ell]\setminus\{i\}}V_j\right)\right)-1 & \text{(since $\bv\notin  V_i+\bigcap_{j\in[\ell]\setminus\{i\}}V_j$)}\\
&=\dim V_{[\ell]}=d.
\end{align*}
Thus, replacing $V_i$ with $V_i'$ preserves the fact that $\dim V_{[\ell]} = d$.
We continue this process until a partition $P_1\sqcup\cdots\sqcup P_s$ with $s\geq 2$ becomes tight. This proves the claim. 

Next, we verify that \cref{eqiv:2} holds for the padded subspaces. For ease of notation, we still denote the padded subspaces by $V_1,\dots,V_\ell$. 
Fix a tight partition  $P_1\sqcup \cdots\sqcup P_s$ of $[\ell]$ with $s\geq 2$, which exists after padding.

Consider arbitrary $i\in [s]$. 
Recall that $V_{P_i}=\bigcap_{j\in P_i}V_j$.
Let $k_i=k-\dim V_{P_i}$, so that $\dim V_{P_i}=k-k_i$.
Let $W_i$ be a complement of $V_{P_i}$ in $\F_q^n$. 
For each $j\in P_i$, 
%let $T_j$ be a complement of $V_{P_i}\subseteq V_j$ in $V_j$.
let $T_j:=V_j\cap W_i$.
For $j\in P_i$,
as $\F_q^n=V_{P_i}\oplus W_i$ and $V_j\supseteq V_{P_i}$, we see that $V_j=V_{P_i}\oplus T_j$ and $\dim T_j = \dim V_j - \dim V_{P_i}\leq k-(k-k_i)=k_i$.


Consider an arbitrary partition $Q_1\sqcup \cdots\sqcup Q_t$ of $P_i$. As $P_1\sqcup\cdots \sqcup P_{i-1}\sqcup Q_1\sqcup \cdots \sqcup Q_t\sqcup P_{i+1}\sqcup\cdots\sqcup P_s$ is a partition of $[\ell]$, by \cref{eqiv:3}, we have
\begin{equation}\label{eq:partition-PQ}
\sum_{j=1}^t \dim V_{Q_j}+\sum_{j\in [s]\setminus \{i\}}\dim V_{P_j} \leq (s+t-2)k+d.
\end{equation}
As the partition $P_1\sqcup\cdots\sqcup P_s$ of $[\ell]$ is tight, we also have
\begin{equation}\label{eq:partition-P}
\sum_{j=1}^s\dim V_{P_j} = (s-1)k+d.
\end{equation}
Define $T_\Omega=\bigcap_{j\in\Omega} T_j$ for nonempty $\Omega\subseteq P_i$.
Then
\begin{align*}
\sum_{j=1}^t \dim T_{Q_j} &=\left(\sum_{j=1}^t \dim V_{Q_j}\right)-t\cdot \dim V_{P_i}\\
&=\left(\sum_{j=1}^t \dim V_{Q_j}+\sum_{j\in [s]\setminus \{i\}}\dim V_{P_j}\right)- \sum_{j=1}^s\dim V_{P_j}-(t-1)\dim V_{P_i}\\
&\stackrel{\eqref{eq:partition-PQ},\eqref{eq:partition-P}}{\leq} (s+t-2)k+d-((s-1)k+d)-(t-1)(k-k_i)=(t-1)k_i.
\end{align*}
This shows that \cref{eqiv:3} holds for the subspaces $(T_j)_{j\in P_i}$ indexed by $P_i$ (instead of by $[\ell]$), where the parameter $d$ is set to zero.

Let $\ell_i=|P_i|$. Note that $\ell_i<\ell$ since $s\geq 2$.
Identifying $P_i$ with $[\ell_i]$ and applying the induction hypothesis, we see that there exist integers $\delta_j\geq 0$ for all $j\in P_i$ such that
\begin{equation}\label{eq:sum-of-delta}
\sum_{j\in P_i}\delta_j=k_i
\end{equation}
and for all nonempty $\Omega\subseteq P_i$,
\begin{equation}\label{eq:T-omega}
\dim T_{\Omega}\leq k_i-\sum_{j\in \Omega} \delta_j.
\end{equation}
As $i\in [s]$ is arbitrary, we can perform the above procedure for $i=1,\dots,s$, which yields $\delta_j$ for all $j\in [\ell]$.
We now verify that $\delta_1,\dots,\delta_{\ell}$ satisfy \cref{eqiv:2}. First, observe that
\[
\sum_{i=1}^\ell \delta_i
=\sum_{i=1}^s\sum_{j\in P_i}\delta_j
\stackrel{\eqref{eq:sum-of-delta}}{=}\sum_{i=1}^s k_i
=\sum_{i=1}^s \bigg(k-\dim V_{P_i}\bigg)\stackrel{\eqref{eq:partition-P}}{=}k-d.
\]
So it remains to verify \eqref{eq:cond2} for all nonempty $\Omega\subseteq [\ell]$. 
Fix nonempty $\Omega\subseteq [\ell]$. Let $I=\{i\in [s]: P_i\cap \Omega \neq \emptyset\}$.
Consider $i\in I$. Recall that for $j\in P_i$, the complement $T_j$ of $V_{P_i}$ in $V_j$ is chosen to be $V_j\cap W_i$, where $W_i$ is a complement of $V_{P_i}$ in $\F_q^n$. It follows that for any $J\subseteq P_i$,
\begin{equation}\label{eq:intersection-of-Vi}
\bigcap_{j\in J} V_j=V_{P_i}\oplus \left(\bigcap_{j\in J} T_j\right)=V_{P_i}\oplus T_{J}.
\end{equation}
Then, we have
\begin{align*}
&\dim\left(\bigcap_{i\in \Omega}V_i\right)
=\dim\left(\bigcap_{i\in I}\bigcap_{j\in P_i\cap \Omega} V_j\right)
\stackrel{\eqref{eq:intersection-of-Vi}}{=}\dim\left(\bigcap_{i\in I}\left(V_{P_i}\oplus T_{P_i\cap\Omega}\right)\right)\\
&\leq\dim\left(\bigcap_{i\in I}V_{P_i}\right)+\sum_{i\in I}\dim\left(T_{P_i\cap\Omega}\right) & \text{(by \cref{lm:dim})}\\
&\stackrel{\eqref{eq:T-omega}}{\leq} \dim\left(\bigcap_{i\in I}V_{P_i}\right)+\sum_{i\in I}\left(k_i-\sum_{j\in P_i\cap\Omega} \delta_j\right)\\
&=\dim\left(\bigcap_{i\in I}V_{P_i}\right)-\sum_{i\in I} \dim V_{P_i}+k|I|-\sum_{j\in \Omega}\delta_j & \text{(as  $\dim V_{P_i}=k-k_i$)}\\
&=\dim\left(\bigcap_{i\in I}V_{P_i}\right)+\sum_{i\in [s]\setminus I}\dim V_{P_i}-\sum_{i\in [s]}\dim V_{P_i}+k|I|-\sum_{j\in \Omega}\delta_j\\
&\stackrel{\eqref{eq:partition-P}}{=}\dim\left(\bigcap_{i\in I}V_{P_i}\right)+\sum_{i\in [s]\setminus I}\dim V_{P_i}-((s-1)k+d)+k|I|-\sum_{j\in \Omega}\delta_j\\
& \leq (s-|I|)k+d-((s-1)k+d)+k|I|-\sum_{j\in \Omega}\delta_j\\ &=k-\sum_{j\in \Omega}\delta_j,
\end{align*}
where the second last step follows by applying \eqref{eq:cond3} to the partition of $[\ell]$ consisting of the set $\bigcup_{i\in I}P_i$ and the sets $P_j$ for $j\in [s]\setminus I$. So \cref{eqiv:2} holds, which concludes the proof.
\end{proof}

\end{document}
